\chapter{The Fundamental Group and Coverings}

\begin{center}\itshape
A rubber band around a missing point cannot shrink to a point without crossing the gap. A rubber band on a sphere can. How can we turn that difference into a calculation? Keep track of loops, and regard two loops as the same when one can be deformed into the other. Going around a circle twice gives the integer $2$. Going around the two holes of a figure eight gives words whose order matters.
\end{center}

Henri Poincar\'e wanted tools that could tell manifolds apart. In his 1895 memoir \emph{Analysis situs}, he introduced the fundamental group and developed the cycle counts now called Betti numbers. The two tools carry different information. His fifth supplement, published in 1904, described a three-dimensional manifold with the homology of $S^3$ but with a fundamental group of order $120$. The cycle counts had missed something that loops could detect.

Poincar\'e then asked whether a closed simply connected three-manifold must be a sphere. That question belongs to a much deeper theory than the one developed here. Grigori Perelman's three preprints of 2002--2003 completed Richard Hamilton's Ricci-flow programme and led to its solution. Perelman declined the Fields Medal in 2006 and the Clay Millennium Prize in 2010. We quote that history; the proofs below concern loops, coverings and graphs.

The first useful covering is the real line wrapped around a circle. A loop downstairs lifts to a path upstairs, and the endpoint records its number of turns. More complicated coverings record subgroups instead of integers. This gives a way to draw algebra: a subgroup of a free group becomes a graph.

Maps are continuous throughout, and $I=[0,1]$. Our examples carry metrics, but the lifting arguments use the compact parameter interval and square, not a metric on the target space.

\section{Loops and the fundamental group}

\begin{definition}
A \emph{path} in $X$ is a map $\gamma:I\to X$. It is a \emph{loop at $x_0$} when $\gamma(0)=\gamma(1)=x_0$. Paths $\gamma_0,\gamma_1$ with the same endpoints are \emph{homotopic}, written $\gamma_0\simeq\gamma_1$, if a map $H:I\times I\to X$ satisfies
\[
H(t,0)=\gamma_0(t),\quad H(t,1)=\gamma_1(t),\quad
H(0,s)=\gamma_0(0),\quad H(1,s)=\gamma_0(1).
\]
Thus the deformation keeps both endpoints fixed. When $\gamma(1)=\delta(0)$, their \emph{product} $\gamma\cdot\delta$ follows $\gamma(2t)$ for $t\le1/2$ and $\delta(2t-1)$ for $t\ge1/2$. Write $\bar\gamma(t)=\gamma(1-t)$ for the reversed path and $c_x$ for the constant path at $x$.
\end{definition}

The parameter $t$ says where you are on a path; $s$ says which stage of the deformation you have reached. Reversing $s$ reverses a homotopy. Joining two homotopies at $s=1/2$ proves transitivity, and a constant homotopy proves reflexivity. Homotopy is therefore an equivalence relation. Its classes are written $[\gamma]$.

\begin{lemma}\label{lem:convexhomotopy}
\begin{enumerate}
\item[(a)] Paths with the same endpoints in a convex subset of $\R^n$ are homotopic through $H(t,s)=(1-s)\gamma_0(t)+s\gamma_1(t)$.
\item[(b)] If $\gamma_0\simeq\gamma_1$ in $X$ and $f:X\to Y$, then $f\circ\gamma_0\simeq f\circ\gamma_1$ in $Y$.
\item[(c)] If $\varphi:I\to I$ is continuous, $\varphi(0)=0$ and $\varphi(1)=1$, then $\gamma\circ\varphi\simeq\gamma$.
\end{enumerate}
\end{lemma}

\begin{proof}
In (a), convexity keeps every segment inside the set, and equal endpoints keep the two sides fixed. In (b), use $f\circ H$. For (c), join $\varphi$ to the identity by straight segments in $I$, then compose with $\gamma$. Monotonicity of $\varphi$ is unnecessary: a path may pause or retrace part of its route without changing its class.
\end{proof}

\begin{proposition}\label{prop:pi1group}
For paths whose indicated products exist:
\begin{enumerate}
\item[(a)] $\gamma\simeq\gamma'$ and $\delta\simeq\delta'$ imply $\gamma\cdot\delta\simeq\gamma'\cdot\delta'$;
\item[(b)] $(\gamma\cdot\delta)\cdot\varepsilon\simeq\gamma\cdot(\delta\cdot\varepsilon)$;
\item[(c)] $c_{\gamma(0)}\cdot\gamma\simeq\gamma\simeq\gamma\cdot c_{\gamma(1)}$;
\item[(d)] $\gamma\cdot\bar\gamma\simeq c_{\gamma(0)}$.
\end{enumerate}
Consequently loop classes at $x_0$ form the \emph{fundamental group} $\pi_1(X,x_0)$, with product $[\gamma][\delta]=[\gamma\cdot\delta]$, identity $[c_{x_0}]$ and inverse $[\bar\gamma]$.
\end{proposition}

\begin{proof}
For (a), put the homotopies side by side in the $t$ direction. Their common edge agrees because the joining endpoint stays fixed. For (b), the left path equals the right path composed with the piecewise linear map taking $0,1/4,1/2,1$ to $0,1/2,3/4,1$. Apply the lemma. For (c), the reparametrization is zero until $1/2$, then $2t-1$; the other order uses a final pause. For (d), use
\[
H(t,s)=
\begin{cases}
\gamma(2t(1-s)),&0\le t\le1/2,\\
\gamma((2-2t)(1-s)),&1/2\le t\le1.
\end{cases}
\]
At stage $s$ the path travels only as far as $\gamma(1-s)$ before returning. At $s=1$ it stays at the starting point. These identities give the group laws on classes.
\end{proof}

For any finite concatenation, changing the parentheses or the time spent on each path changes only its parametrization. In particular we may assign equal time to all its pieces. This will let us subdivide products of loops without keeping track of their original speeds.

A path $h$ from $x_0$ to $x_1$ changes basepoints by
\[
\beta_h:\pi_1(X,x_0)\longrightarrow\pi_1(X,x_1),\qquad
[\gamma]\longmapsto[\bar h\gamma h].
\]
The adjacent $h\bar h$ cancels when products are formed, so this is a homomorphism. Its inverse is $\beta_{\bar h}$. In a path-connected space we can therefore write $\pi_1(X)$ when only its isomorphism type matters. A space is \emph{simply connected} if it is path-connected and this group is trivial. Every nonempty convex subset of $\R^n$ is simply connected by the lemma.

A map $f:X\to Y$ induces $f_*[\gamma]=[f\circ\gamma]$ from $\pi_1(X,x_0)$ to $\pi_1(Y,f(x_0))$. Composition preserves homotopies and concatenations, so this is a homomorphism. Also $(gf)_*=g_*f_*$ and $\mathrm{id}_*=\mathrm{id}$. A homeomorphism consequently induces an isomorphism.

\begin{proposition}\label{prop:homotopyinv}
Suppose $F:X\times I\to Y$ joins $f$ to $g$, and put $h(s)=F(x_0,s)$. Then $g_*=\beta_h f_*$. If $f$ is a \emph{homotopy equivalence}, meaning that some $g:Y\to X$ has $gf\simeq\mathrm{id}_X$ and $fg\simeq\mathrm{id}_Y$, then $f_*$ is an isomorphism.
\end{proposition}

\begin{proof}
For a loop $\gamma$, map the square by $(t,s)\mapsto F(\gamma(t),s)$. In the square the bottom edge can be deformed, with endpoints fixed, into the route up the left side, across the top and down the right side. The lemma carries that deformation into $Y$, giving $f\gamma\simeq h(g\gamma)\bar h$. Moving $h$ to the other side gives the formula.

For a homotopy inverse, $g_*f_*=(gf)_*$ is an isomorphism, after the basepoint change just proved. Hence $f_*$ is injective and $g_*$ is surjective. Apply the same argument to $fg$ at $f(x_0)$: its first map is this same $g_*$, so $g_*$ is injective as well. Now $g_*$ and $g_*f_*$ are bijective, which makes $f_*$ bijective.
\end{proof}

A \emph{deformation retraction} of $X$ onto $A\subset X$ is a homotopy from the identity to a map into $A$ that fixes every point of $A$ throughout. The inclusion and final map are homotopy inverses, so $\pi_1(A,a)\cong\pi_1(X,a)$. For example,
\[
F(x,s)=(1-s)x+s\frac{x}{|x|}
\]
retracts $\R^2\setminus\{0\}$ onto its unit circle: each point moves along its positive radial ray. A solid torus retracts onto its central circle by contracting each cross-sectional disc.

\section{Covering spaces and the circle}

\begin{definition}
A map $p:\tilde X\to X$ is a \emph{covering} if each point of $X$ has an open neighbourhood $U$ whose inverse image is a disjoint union of open sets, each mapped homeomorphically onto $U$. Such a $U$ is \emph{evenly covered}, and the open sets upstairs are its \emph{sheets}. A \emph{lift} of $f:Y\to X$ is a map $\tilde f:Y\to\tilde X$ satisfying $p\tilde f=f$.
\end{definition}

This definition allows an empty collection of sheets on a component outside the image. Our based, connected coverings lie over the component containing their chosen basepoint; there path lifting will show that every fibre is nonempty.

\begin{example}\label{ex:coverings}
(a) The map $p(t)=e^{2\pi it}$ covers $S^1$ by $\R$. A proper open arc lifts to disjoint translates of an interval, one for each integer. Figure~\ref{fig:helix} draws those intervals along a helix.
(b) For a nonzero integer $n$, $z\mapsto z^n$ covers $S^1$ with $|n|$ sheets: choose arcs short enough that their arguments have $|n|$ separate inverse branches. The constant map for $n=0$ is not a covering.
(c) Small boxes in $T^2=\R^2/\Z^2$ lift to their integer translates in $\R^2$.
(d) For $n\ge0$, $S^n\to\R P^n$ identifies antipodal points. An open set contained in an open hemisphere and its antipode give the two sheets over its image. For $n=0$ this is the two-point space covering a point.
\end{example}

\begin{figure}[ht]
\centering
\fig{ch05-helix}
\caption{The covering $\R\to S^1$ drawn as a helix over the circle. A loop that goes twice around the circle lifts to a path that climbs two turns of the helix: the
height it gains is the number of turns, and no homotopy of the loop can change it.}
\label{fig:helix}
\end{figure}

\begin{lemma}[uniqueness of lifts]\label{lem:uniquelift}
If $Y$ is connected, two lifts of the same map $f:Y\to X$ through a covering $p:\tilde X\to X$ that agree at one point agree everywhere.
\end{lemma}

\begin{proof}
Near $f(y)$ choose an evenly covered $U$. The two lifted values lie in sheets $U_1,U_2$. On the common inverse image of those sheets near $y$, equal sheets force the lifts to agree, since $p$ is injective there. Different sheets force them to differ, since sheets are disjoint. Thus both the agreement set and its complement are open. Connectedness and one point of agreement finish the proof.
\end{proof}

\begin{theorem}[homotopy lifting]\label{thm:homotopylifting}
For a covering $p:\tilde X\to X$, a map $F:I\times I\to X$ and $\tilde x$ with $p(\tilde x)=F(0,0)$, there is exactly one lift $\tilde F$ with $\tilde F(0,0)=\tilde x$. In particular every path has a unique lift from any specified point above its start.
\end{theorem}

\begin{proof}
The uniqueness lemma applies because the square is connected. For existence, the sets $F^{-1}(U)$ for evenly covered $U$ cover the compact square. Their Lebesgue number (Volume~7, Theorem~4.8) gives a finite grid of small closed squares, each mapped into one such $U$.

Start at the lower left square and work along each row, from bottom to top. On the first square choose the sheet containing $\tilde x$ and use its inverse branch. Every later square meets the part already treated in its left edge, its bottom edge, or their connected union. The lift on this intersection lies in a single sheet over the new $U$: its image is connected and sheets are disjoint open sets. That sheet's inverse branch extends the lift to the whole square and agrees with the earlier lift on the overlap. The finitely many closed pieces glue to a continuous map. Taking $F(t,s)=\gamma(t)$ gives path lifting.
\end{proof}

\begin{corollary}\label{cor:liftendpoints}
Lifts from the same starting point of homotopic paths have the same endpoint and are homotopic with endpoints fixed. Hence $p_*:\pi_1(\tilde X,\tilde x_0)\to\pi_1(X,p(\tilde x_0))$ is injective.
\end{corollary}

\begin{proof}
Lift a homotopy square. Its left and right edges project to constant paths, whose lifts are constant by uniqueness. Its bottom and top are the prescribed path lifts. This is the required fixed-endpoint homotopy. If a loop upstairs projects to a nullhomotopic loop, lift a contraction; its final constant edge lifts to a constant. The original loop upstairs therefore contracts too.
\end{proof}

\begin{theorem}\label{thm:pi1circle}
The map $\pi_1(S^1,1)\to\Z$ sending $[\gamma]$ to $\tilde\gamma(1)$, where $\tilde\gamma$ lifts $\gamma$ through $p(t)=e^{2\pi it}$ from $0$, is an isomorphism. It sends $[\omega_n]$ to $n$ for $\omega_n(t)=e^{2\pi int}$.
\end{theorem}

\begin{proof}
The lifted endpoint is an integer because its exponential is $1$, and it depends only on the class by the corollary. If $\tilde\gamma(1)=m$, the lift of $\gamma\delta$ is $\tilde\gamma$ followed by $m+\tilde\delta$. Its endpoint is $m+\tilde\delta(1)$, so endpoints add. The path $nt$ lifts $\omega_n$, giving surjectivity. Finally, a lift ending at zero is a loop in $\R$. Contract it there by straight segments and project the contraction. This proves injectivity.
\end{proof}

The lift is a continuous argument divided by $2\pi$, as in Volume~7, Lemma~6.4. Its endpoint is the winding number. After radial retraction, loops in $\C^*$ with the \emph{same basepoint} are homotopic exactly when their winding numbers agree.

For any covering the \emph{fibre} over $x_0$ is $p^{-1}(x_0)$. A loop class $g$ acts on it on the right: lift a representative from $\tilde x$ and call its endpoint $\tilde x\cdot g$. The corollary makes the result independent of the representative. Concatenating lifts gives $(\tilde x\cdot g)\cdot h=\tilde x\cdot(gh)$.

\begin{proposition}\label{prop:fibre}
Suppose $\tilde X$ is path-connected, $p(\tilde x_0)=x_0$, and $H=p_*\pi_1(\tilde X,\tilde x_0)$. The action on the fibre over $x_0$ is transitive and its stabilizer at $\tilde x_0$ is $H$. Consequently $\tilde x_0\cdot g\mapsto Hg$ identifies that fibre with the right cosets of $H$ in $\pi_1(X,x_0)$. Its cardinality is the index $[\pi_1(X,x_0):H]$.
\end{proposition}

\begin{proof}
Join $\tilde x_0$ to another point of the fibre by a path upstairs. Its projection is a loop, and lifting it recovers the chosen path, proving transitivity. A class fixes $\tilde x_0$ precisely when its lift is a loop. Such classes are exactly the image under $p_*$ of loop classes upstairs, namely $H$. Finally $\tilde x_0g=\tilde x_0h$ is equivalent to $gh^{-1}\in H$, or $Hg=Hh$.
\end{proof}

Path lifting identifies fibres over points joined by a path, so their cardinalities agree on each path component. Over a path-connected base this cardinality is the number of sheets.

From here on our spaces are path-connected and \emph{locally path-connected}: every neighbourhood contains an open path-connected neighbourhood of its point. Manifolds and the graphs used below have this property.

\begin{theorem}[lifting criterion]\label{thm:liftcrit}
Let $p:(\tilde X,\tilde x_0)\to(X,x_0)$ be a covering and $f:(Y,y_0)\to(X,x_0)$ a map, with $Y$ path-connected and locally path-connected. A lift taking $y_0$ to $\tilde x_0$ exists exactly when
\[
f_*\pi_1(Y,y_0)\subset p_*\pi_1(\tilde X,\tilde x_0).
\]
When it exists, the lift is unique.
\end{theorem}

\begin{proof}
A lift gives $f_*=p_*\tilde f_*$, proving necessity. For the converse, join $y_0$ to $y$ by a path $\eta$ and define $\tilde f(y)$ as the endpoint of the lift of $f\eta$ from $\tilde x_0$. If $\eta'$ is another choice, the hypothesis says that $f\eta\,\overline{f\eta'}$ has a closed lift. Indeed its class is the projection of a loop upstairs, and the endpoint corollary applies. Reversing the second part of that closed lift shows that $f\eta$ and $f\eta'$ lift to the same endpoint. Thus the definition is well defined.

Choose an evenly covered $U$ about $f(y)$ and the sheet $\tilde U$ containing $\tilde f(y)$. Local path-connectedness gives an open path-connected $W$ about $y$ with $f(W)\subset U$. For $y'\in W$, extend $\eta$ by a path inside $W$. Its image lifts using the inverse branch on $\tilde U$, so on $W$ we have $\tilde f=(p|_{\tilde U})^{-1}f$. This proves continuity. The uniqueness lemma applies to $Y$.
\end{proof}

\section{Deck transformations and the Galois correspondence}

\begin{definition}
A \emph{deck transformation} is a homeomorphism $\tau:\tilde X\to\tilde X$ with $p\tau=p$. These form the group $\operatorname{Deck}(p)$. Coverings $p_i:\tilde X_i\to X$ are \emph{isomorphic} if a homeomorphism $\phi:\tilde X_1\to\tilde X_2$ satisfies $p_2\phi=p_1$. A covering with path-connected total space is \emph{normal} if its deck group acts transitively on every fibre, and \emph{universal} if its total space is simply connected.
\end{definition}

For $\R\to S^1$, deck transformations are integer translations. For $z\mapsto z^n$ on $S^1$, they are rotations by $n$th roots of unity. In general a deck transformation is determined by the image of one point: it and any other candidate are lifts of $p$, so uniqueness applies. It also carries each lifted path to the lift with the new starting point. Hence
\[
\tau(\tilde x\cdot g)=\tau(\tilde x)\cdot g.
\]

\begin{theorem}\label{thm:deck}
For a covering with path-connected $\tilde X$, write $G=\pi_1(X,x_0)$ and $H=p_*\pi_1(\tilde X,\tilde x_0)$.
\begin{enumerate}
\item[(a)] The relation $\tilde x_1=\tilde x_0\cdot g$ gives $p_*\pi_1(\tilde X,\tilde x_1)=g^{-1}Hg$.
\item[(b)] A deck transformation taking $\tilde x_0$ to $\tilde x_0\cdot g$ exists exactly when $g\in N(H)$, the normalizer of $H$.
\item[(c)] $\operatorname{Deck}(p)\cong N(H)/H$. The covering is normal exactly when $H$ is normal in $G$; then its deck group is $G/H$. A universal covering has deck group $G$.
\end{enumerate}
\end{theorem}

\begin{proof}
(a) Lift a representative $\gamma$ of $g$ from $\tilde x_0$ to $\tilde x_1$. Changing basepoints along that lift replaces a projected loop $\lambda$ by $\bar\gamma\lambda\gamma$. This conjugates its class by $g^{-1}$ and gives the asserted subgroup.

(b) Apply the lifting criterion to $p$, asking for a lift $\tau$ taking $\tilde x_0$ to $\tilde x_1$. It exists when $H\subset g^{-1}Hg$. A lift in the opposite direction exists when the reverse inclusion holds. If both hold, their composites fix the specified points and lift $p$, so uniqueness makes both composites identities. They are inverse homeomorphisms. Conversely a deck homeomorphism supplies both lifts. Thus we need equality, which is exactly $g\in N(H)$.

(c) Let $\tau_g$ denote the transformation in (b). The commutation rule above gives
\[
\tau_g\tau_h(\tilde x_0)=\tau_g(\tilde x_0\cdot h)=\tilde x_0\cdot gh,
\]
so $\tau_g\tau_h=\tau_{gh}$. Every deck transformation sends $\tilde x_0$ to some point of its fibre and hence arises this way. The kernel consists of $g$ fixing $\tilde x_0$, namely $H$. Transitivity on the fibre is equivalent to $N(H)=G$. Finally, a simply connected total space has $H=\{1\}$.
\end{proof}

\begin{example}\label{ex:computations}
(a) For $n\ge0$, $\R^n\to T^n=\R^n/\Z^n$ is universal. A deck transformation has $\tau(x)-x\in\Z^n$; continuity and connectedness make this difference constant. Conversely each integer translation is a deck transformation. Thus $\pi_1(T^n)=\Z^n$, including the trivial group for the one-point torus $T^0$.
(b) Once we prove Corollary~\ref{cor:spheres} below, the same method will give $\pi_1(\R P^n)=\Z/2$ for $n\ge2$: the universal covering is $S^n$, and its two deck maps are the identity and the antipodal map. This is a preview of that computation.
(c) For $n\ge1$, the circle covering $z\mapsto z^n$ sends the generator of $\pi_1(S^1)$ to $n$, so $H=n\Z$. It is normal, with deck group $\Z/n$, realized by multiplication by $n$th roots of unity.
\end{example}

The subgroup correspondence resembles Galois theory in Volume~3, Chapter~10: intermediate coverings take the place of intermediate fields, and normal subgroups give the coverings with all possible fibre symmetries. To complete the correspondence we must construct a covering from a subgroup. A local condition makes that possible.

A space is \emph{semilocally simply connected} if every point $x$ has a neighbourhood $U$ such that each loop in $U$ based at $x$ contracts in $X$. The contraction need not stay in $U$. A manifold satisfies this condition by taking a small coordinate ball. The small stars of graphs will give another example.

\begin{theorem}[classification of coverings]\label{thm:galois}
Suppose $X$ is path-connected, locally path-connected and semilocally simply connected. For every $H\le\pi_1(X,x_0)$ there is a covering
$p_H:(X_H,\tilde x_0)\to(X,x_0)$ with path-connected $X_H$ and $p_{H*}\pi_1(X_H,\tilde x_0)=H$. Two based coverings with these properties and the same subgroup have a unique basepoint-preserving isomorphism. The subgroup $\{1\}$ gives a universal covering.
\end{theorem}

\begin{proof}
Call an open set \emph{good} if it is path-connected and all its loops contract in $X$. Good sets form a basis. Indeed, at $x$ choose a neighbourhood with the semilocal property, then a path-connected open subset $W$ inside it and inside any prescribed neighbourhood. A loop in $W$ based at another point can be transported to $x$ along a path in $W$; its transported loop contracts in $X$, so the original does too.

A point upstairs will remember a path from $x_0$, while forgetting precisely the loops in $H$. For paths starting at $x_0$, set
\[
\gamma\sim\gamma'\quad\Longleftrightarrow\quad
\gamma(1)=\gamma'(1)\ \hbox{and}\ [\gamma\bar\gamma']\in H.
\]
Reflexivity, symmetry and transitivity follow from the subgroup laws and cancellation of a reversed path. Let $X_H$ be the set of classes $\langle\gamma\rangle$, and set $p_H\langle\gamma\rangle=\gamma(1)$.

For a good $U$ containing $\gamma(1)$, define
\[
U_{\langle\gamma\rangle}=
\{\langle\gamma\eta\rangle:\eta\text{ is a path in }U\text{ starting at }\gamma(1)\}.
\]
This depends only on the class of $\gamma$: appending the same $\eta$ cancels in $\gamma\eta\overline{\gamma'\eta}$. The endpoint map is a bijection from this set to $U$. Surjectivity uses path-connectedness. For injectivity, two paths $\eta,\eta'$ with the same endpoint differ by a loop in $U$, nullhomotopic in $X$; hence $\gamma\eta\sim\gamma\eta'$.

If $\langle\gamma'\rangle\in U_{\langle\gamma\rangle}$, appending paths in $U$ shows $U_{\langle\gamma'\rangle}=U_{\langle\gamma\rangle}$; use the reverse connecting path for the other inclusion. Thus two such sets over the same $U$ are equal or disjoint. At a point of two basic sets, choose a good $W$ about its endpoint inside their intersection. The corresponding $W_{\langle\gamma'\rangle}$ lies inside both. Our sets therefore form a topology basis. Endpoint maps take smaller basic sets to the corresponding smaller good sets, so $p_H$ restricts to a homeomorphism on each $U_{\langle\gamma\rangle}$. These sets partition $p_H^{-1}(U)$. We have constructed a covering.

Take $\tilde x_0=\langle c_{x_0}\rangle$. For $\gamma$ starting at $x_0$, put $\gamma_t(u)=\gamma(tu)$. The path $t\mapsto\langle\gamma_t\rangle$ is continuous: locally its endpoints stay in a good $U$, where it is exactly the inverse branch over $U$ applied to $\gamma(t)$. The identification follows by joining or reversing the short intervening piece of $\gamma$ and cancelling retracing. This path joins $\tilde x_0$ to $\langle\gamma\rangle$, proving path-connectedness. A loop closes upstairs exactly when $\langle\gamma\rangle=\langle c_{x_0}\rangle$, or $[\gamma]\in H$. Proposition~\ref{prop:fibre} identifies the covering subgroup as $H$.

For uniqueness, apply the lifting criterion in both directions between two based coverings with subgroup $H$. The composites lift the respective projections and fix the chosen points. Uniqueness of lifts makes them identities. When $H$ is trivial, injectivity of $p_{H*}$ also makes $\pi_1(X_H)$ trivial, so $X_H$ is simply connected.
\end{proof}

Moving the chosen point in a fibre conjugates $H$. Consequently unbased connected coverings correspond to conjugacy classes of subgroups. The whole group gives a one-sheeted covering, isomorphic to $X$; the trivial subgroup gives the universal covering. The latter covers every connected covering: its projection lifts through that covering by the criterion. On sufficiently small good sets this lift is a homeomorphism between sheets, and path lifting shows it reaches every point of the connected total space. It is therefore itself a covering.

\section{Free products and the theorem of Seifert and van Kampen}

A circle contributes one generator. What happens when two circles meet at a point? A loop may go around the first, the second, then the first backwards. There is no reason to commute those moves. The algebra that keeps them separate is the free product.

For groups $(G_\alpha)$, a \emph{reduced word} is a finite sequence of nonidentity elements of the factors, with successive letters in different factors. The empty word is allowed. Write $W$ for their set.

\begin{proposition}\label{prop:freeproduct}
The reduced words form a group $\mathop{\ast}_\alpha G_\alpha$, the \emph{free product}. Each factor embeds as words of length at most one. To multiply, concatenate words, multiply adjacent letters from the same factor and delete identities until the word is reduced. Homomorphisms $\varphi_\alpha:G_\alpha\to K$ extend uniquely to a homomorphism $\mathop{\ast}_\alpha G_\alpha\to K$.
\end{proposition}

\begin{proof}
We must justify that reduction gives an associative operation. For $g\in G_\alpha$, let $L_g$ act on $W$ by inserting $g$ on the left and doing the required first reduction. If $g=1$, do nothing. If the first letter is in another factor, prepend $g$. If it is $a\in G_\alpha$, replace it by $ga$, deleting it when $ga=1$.

For $g,h\in G_\alpha$, we have $L_gL_h=L_{gh}$. To check this, first take a word empty or beginning in a different factor: prepending $h$ and then combining $g$ with it gives $gh$, with deletion if needed. For a word $aw'$ beginning in $G_\alpha$, if $ha\ne1$, both procedures replace $a$ by $g(ha)=(gh)a$, deleting an identity. If $ha=1$, the first procedure removes $a$, then prepends $g$ to $w'$, whose first factor is different. The second combines $(gh)a=g$ and gives the same result. Cases with identity letters are included by $L_1=\mathrm{id}$. Thus $L_g$ is a permutation with inverse $L_{g^{-1}}$.

Let $P$ be the group generated by these permutations. Combining adjacent operators from the same factor expresses every member as $L_{g_1}\cdots L_{g_m}$ for a reduced word. Applied to the empty word, this permutation returns $g_1\cdots g_m$. Different reduced words therefore give different permutations. Transporting the group law of $P$ to $W$ proves the asserted multiplication and associativity.

Finally send a reduced word to $\varphi(g_1)\cdots\varphi(g_m)$ in $K$. Each allowed reduction leaves this value unchanged because each $\varphi_\alpha$ is a homomorphism. Thus the map respects products. Its values on the factors force uniqueness. If there are no nontrivial factors, only the empty word remains, giving the trivial group.
\end{proof}

The \emph{free group} $F(S)$ is the free product of one copy of $\Z$ for each $s\in S$. Its elements are words in $s^{\pm1}$ with adjacent inverse pairs cancelled. A map from $S$ into any group extends uniquely to a homomorphism. Write $F_n$ for $n$ generators; $F_0$ is trivial. Volume~8, Chapter~13 realized $F_2$ inside the rotation group in the Banach--Tarski argument.

For maps $i_j:A\to G_j$, the \emph{amalgamated product} $G_1\ast_A G_2$ is $(G_1\ast G_2)/N$, where $N$ is the normal subgroup generated by $i_1(a)i_2(a)^{-1}$. The normal subgroup generated by a set is the smallest normal subgroup containing it, equivalently all finite products of conjugates of its elements and their inverses. The quotient identifies the two copies of every $a$, even when the maps $i_j$ are not injective.

\begin{theorem}[Seifert--van Kampen]\label{thm:vankampen}
Suppose $X=U\cup V$, where $U,V$ are open and $U,V,U\cap V$ are path-connected, with $x_0\in U\cap V$. Let $i_U,i_V$ be the inclusion maps on fundamental groups from $U\cap V$, and $j_U,j_V$ those into $X$. The homomorphism
\[
\Phi:\pi_1(U,x_0)\ast\pi_1(V,x_0)\longrightarrow\pi_1(X,x_0)
\]
extending $j_U,j_V$ is onto. Its kernel is the normal subgroup $N$ generated by $i_U(\omega)i_V(\omega)^{-1}$ for $\omega\in\pi_1(U\cap V,x_0)$. Thus
\[
\pi_1(X)\cong\pi_1(U)\ast_{\pi_1(U\cap V)}\pi_1(V).
\]
\end{theorem}

\begin{proof}
A loop can pass between $U$ and $V$ away from $x_0$. We first connect each passage back to $x_0$ to obtain based loops. Subdivide a loop $f$ into finitely many paths $f_i$ lying in $A_i\in\{U,V\}$, using a Lebesgue number on $I$. At the joining points choose paths $g_i$ from $x_0$ inside $A_i\cap A_{i+1}$; take the two endpoint paths constant. Then
\[
f\simeq(f_1\bar g_1)(g_1f_2\bar g_2)\cdots(g_{m-1}f_m).
\]
Each factor lies in one of the two sets. This proves surjectivity. The generators of $N$ project to the identity because the two inclusions of $U\cap V$ into $X$ are the same, so $N\subset\ker\Phi$.

Put $Q=(\pi_1(U)\ast\pi_1(V))/N$. Combining adjacent loops in one set does not change their product in $Q$. A loop in $U\cap V$ can also be regarded as a loop in either set without changing its value in $Q$. We show that these facts make any two factorizations of the same loop class equal in $Q$.

Represent the two factorizations by loops $f_1,\ldots,f_k$ and $f'_1,\ldots,f'_l$, assigning equal time to the pieces as justified earlier. Replace an empty factorization by a constant loop, so $k,l\ge1$. Let $F$ be a based homotopy between the concatenations. Only the left and right edges of its square are necessarily constant. Along the bottom and top, $F$ returns to $x_0$ at the factor joins $i/k$ and $j/l$.

Take a common vertical subdivision through all those joins, then refine both directions until every rectangle $R$ maps into a chosen $A(R)\in\{U,V\}$. At each grid vertex $v$, choose a path $g_v$ from $x_0$ to $F(v)$ inside all incident $A(R)$ and, on the bottom or top, inside the sets containing the corresponding factor loops. This is possible: every such intersection is $U$, $V$ or $U\cap V$, is path-connected, and contains both endpoints. Choose $g_v$ constant whenever $F(v)=x_0$, in particular on the left and right sides and at factor joins.

For an oriented grid edge from $v$ to $w$, let
\[
L(v,w)=g_v\cdot F|_{[v,w]}\cdot\bar g_w.
\]
It lies in every incident rectangle's chosen set and in the appropriate boundary factor set. If these choices differ, the loop lies in $U\cap V$ and has the same value in $Q$ with either interpretation.

In one rectangle, the products associated with its bottom-right and left-top routes agree in $\pi_1(A(R))$. Cancel the middle $g$ paths; the remaining two routes are homotopic through $F$ on the rectangle. Consequently swapping these routes does not change a grid path's product in $Q$. Successively swapping across the finitely many rectangles changes the route along the bottom and right sides into the route along the left and top sides.

The initial product equals the first factorization: the right-side loops are constant, and the bottom-edge pieces within a factor combine in that factor's set; their connector paths cancel and the factor endpoints have constant connectors. The final product equals the second factorization for the same reason. The two words therefore agree in $Q$. Apply this to a word in $\ker\Phi$ and the constant factorization. Its class in $Q$ is the identity, so it belongs to $N$.
\end{proof}

\begin{corollary}\label{cor:spheres}
For $n\ge2$, $S^n$ is simply connected.
\end{corollary}

\begin{proof}
Remove the north pole to obtain $U$ and the south pole to obtain $V$. Stereographic projection from Exercise~\ref{ex:stereographic} identifies each with the convex space $\R^n$. Their intersection corresponds to $\R^n\setminus\{0\}$, path-connected when $n\ge2$. Van Kampen makes $\pi_1(S^n)$ a quotient of two trivial groups, hence trivial.
\end{proof}

For $S^1$, that intersection has two components, so the theorem does not apply; its group is $\Z$. The two-point space $S^0$ is not path-connected.

Another useful gluing adds a disc. Given $\varphi:S^1\to X$, write $X\cup_\varphi D^2$ for the quotient that identifies each boundary point $z$ with $\varphi(z)$. The disc lets its boundary loop contract, so we expect to kill that loop and all its conjugates.

\begin{proposition}\label{prop:attachdisc}
For path-connected $X$,
\[
\pi_1(X\cup_\varphi D^2)\cong\pi_1(X)/\langle\langle[\varphi]\rangle\rangle.
\]
Here $[\varphi]$ is transported to the chosen basepoint along any path; the double brackets mean its normal closure.
\end{proposition}

\begin{proof}
In the attached space $Y$, take $U$ to be the complement of the disc's centre and $V$ its open interior. These are open. Pushing points of the punctured disc radially to its boundary retracts $U$ onto $X$. The set $V$ is simply connected, while $U\cap V$ is a punctured disc with group $\Z$. At basepoint $1/2$ in the disc, its generating circle retracts to $\varphi$, with a basepoint change along the radius. Van Kampen therefore kills exactly the normal closure of this class. A different path to the original basepoint conjugates the class and leaves its normal closure unchanged.
\end{proof}

Before using this rule for surfaces, compute the group of a bouquet $R_m$ of $m$ circles. For $m=0$ it is a point. For $m\ge1$, remove the midpoint of the last circle to form $U$; it retracts onto $R_{m-1}$. Let $V$ be the last circle together with short open stubs of the other circles. It retracts onto that circle, while $U\cap V$ is a contractible star. Van Kampen gives $\pi_1(R_m)=\pi_1(R_{m-1})\ast\Z$. Induction proves $\pi_1(R_m)=F_m$, with the individual circles as its free generators.

\begin{example}\label{ex:surfaces}
For $g\ge1$, identify the sides of a $4g$-gon in the order
\[
a_1b_1a_1^{-1}b_1^{-1}\cdots a_gb_ga_g^{-1}b_g^{-1}.
\]
The result is the closed orientable surface $\Sigma_g$. Its vertices become one point, its edges a bouquet of $2g$ circles, and its interior an attached disc. The bouquet computation and the proposition give
\[
\pi_1(\Sigma_g)\cong
\langle a_1,b_1,\ldots,a_g,b_g\mid [a_1,b_1]\cdots[a_g,b_g]\rangle,
\qquad [a,b]=aba^{-1}b^{-1}.
\]
This presentation means the free group modulo the normal closure of the displayed word. Its \emph{abelianization}, the quotient by all commutators, is $\Z^{2g}$: after commuting the generators, the relator imposes no further relation, and the free abelian universal property gives the identification. Different genera have different abelianization ranks, hence different fundamental groups. They cannot be homeomorphic or homotopy equivalent. For $g=1$ the relation says $ab=ba$, recovering $\Z^2$. For $g=0$ the surface is $S^2$ and its group is trivial, consistent with no generators.

The projective plane has one circular edge $a$ and a disc whose boundary reads $a^2$, because antipodal boundary points are identified. Its group is therefore $\langle a\mid a^2\rangle=\Z/2$, agreeing with the sphere covering.
\end{example}

\section{Graphs and free groups}\label{sec:graphs}

A \emph{graph} consists of vertices joined by edges, each a copy of $[0,1]$ with its endpoints attached to vertices. An edge may return to its own vertex, and two edges may have the same endpoints. Give every edge length $1$ and use the resulting path metric. A graph is \emph{locally finite} when only finitely many edge ends meet each vertex; a loop edge contributes two ends. Small neighbourhoods of interior points are intervals, and small vertex neighbourhoods are stars. A star contracts along its arms to its centre. Thus these graphs are locally path-connected and semilocally simply connected, even if the number of arms is infinite.

A connected graph here has at least one vertex and an edge path between every pair of vertices. A \emph{tree} is a connected graph containing no embedded circle. Between any two vertices of a tree there is exactly one edge path without backtracking: two different such paths would produce a circle. A \emph{spanning tree} contains all the vertices of its graph.

\begin{lemma}\label{lem:trees}
\begin{enumerate}
\item[(a)] If a vertex $v$ has just one incident edge end, on a nonloop edge $e$, then $Y$ deformation retracts onto $Y\setminus(\{v\}\cup e^\circ)$, where $e^\circ$ is the edge interior.
\item[(b)] A finite tree with $V$ vertices has $V-1$ edges and retracts onto any chosen vertex. A finite tree with one extra edge is homotopy equivalent to a circle.
\item[(c)] A spanning tree exists in every finite connected graph.
\end{enumerate}
\end{lemma}

\begin{proof}
For (a), slide the leaf edge linearly into its other endpoint, fixing the rest of the graph. For (b), a longest path without backtracking in a finite tree ends at a leaf. Unless the tree is a single vertex, choose a leaf other than the prescribed vertex, remove it by (a), and induct. Each removal deletes one vertex and one edge, proving the count. A tree plus an extra edge has exactly one circle. Remove leaves until only that circle remains; equivalently, the paths attached to the circle are finite trees and contract onto their attachment points. A loop edge counts as a circle in this argument. For (c), delete an edge from a circle whenever a circle remains. The alternate route around that circle preserves connectedness. Finitely many deletions leave a connected graph with no circle.
\end{proof}

\begin{theorem}\label{thm:wedge}
Let $X$ be a finite connected graph with $V$ vertices, $E$ edges, spanning tree $T$ and vertex $x_0$. For each oriented edge $e:u\to w$ outside $T$, let $\gamma_e$ go in $T$ from $x_0$ to $u$, traverse $e$, and return in $T$ from $w$. Then $\pi_1(X,x_0)$ is freely generated by the $[\gamma_e]$ and has rank $E-V+1$. A bouquet of $n$ circles therefore has group $F_n$.
\end{theorem}

\begin{proof}
Induct on the number of edges outside $T$. With none, the graph is a contractible finite tree. Otherwise choose $e\notin T$ and remove its midpoint to obtain $U$. Let $V'$ be the open edge $e^\circ$ together with the open $1/3$-neighbourhood of $T$ in $X$. The set $U$ retracts onto $X'=X\setminus e^\circ$. The set $V'$ retracts onto $T\cup e$, a tree with one extra edge; its group is generated by $[\gamma_e]$. Indeed its unique circle is traversed once by that loop, with tree paths before and after. The intersection $U\cap V'$ retracts onto $T$: all its other pieces are short edge stubs, including the two halves of $e$ without the midpoint. All three sets are path-connected. Van Kampen gives
\[
\pi_1(X,x_0)\cong\pi_1(X',x_0)\ast\langle[\gamma_e]\rangle.
\]
Induction gives the claimed free basis. There are $E-(V-1)$ edges outside the spanning tree.
\end{proof}

\begin{figure}[ht]
\centering
\fig{ch05-coverings}
\caption{The figure eight $R_2$, with loops $a$ (coral) and $b$ (charcoal), and its three connected two-sheeted coverings. Each vertex has exactly one outgoing and one incoming edge of each
colour; reading edge labels along loops at the lower vertex gives the subgroup of index $2$: $\langle a,\,b^2,\,bab^{-1}\rangle$, $\langle b,\,a^2,\,aba^{-1}\rangle$ and $\langle a^2,\,ab,\,b^2\rangle$, each free of rank $3$.}
\label{fig:coverings}
\end{figure}

For $H\le F_n=\pi_1(R_n)$ we can build the covering without abstract path classes. Its \emph{Schreier coset graph} has vertices $Hg$, and an edge labelled $a_i$ from $Hg$ to $Hga_i$. Right multiplication permutes cosets, so every vertex has exactly one incoming and one outgoing edge of each label. Mapping each labelled edge to its corresponding circle in $R_n$ gives a covering; a small star downstairs has one identical star over it at each vertex. The graph is connected because every coset has a word representative. Reading a word $w$ from $H$ ends at $Hw$. Thus closed words are exactly $H$, and the covering subgroup is $H$.

\begin{figure}[ht]
\centering
\fig{v9-ch05-three-sheets}
\caption{A three-sheeted cover of the figure eight. The $a$-edges cycle $1\to2\to3\to1$, while the $b$-edges exchange $1,2$ and fix $3$. Each label has one incoming and one outgoing edge at every vertex. A symmetry of the covering must preserve both permutations; only the identity does. Exercise~\ref{ex:v9-three-sheet-cover} computes its subgroup.}
\label{fig:v9-three-sheet-cover}
\end{figure}

\begin{figure}[ht]
\centering
\fig{v9-ch05-shrinking-circles}
\caption{The circles with centres $(1/n,0)$ and radii $1/n$ all pass through $0$. Only the first few are drawn. Every neighbourhood of $0$ contains whole circles, however small it is; each such circle still detects a nontrivial loop. Exercise~\ref{ex:v9-shrinking-circles} uses this to rule out a simply connected covering.}
\label{fig:v9-shrinking-circles}
\end{figure}

\begin{theorem}[Nielsen--Schreier]\label{thm:nielsenschreier}
Every subgroup $H$ of $F_n$ is free. If its index $k$ is finite, its rank is
\[
1+k(n-1).
\]
\end{theorem}

\begin{proof}
The coset graph $X_H$ covers $R_n$, and injectivity of the induced map identifies $\pi_1(X_H)$ with $H$. If the index is $k<\infty$, there are $k$ vertices and $nk$ edges. The graph theorem gives rank $nk-k+1$. When $n=0$, the only subgroup has $k=1$ and rank zero, as the formula says.

For infinite index, $X_H$ is countable and locally finite. Root it at $H$. At every other vertex choose one edge to a vertex one step nearer the root. Following chosen edges reaches the root, so they connect all vertices. They have no circle: a vertex farthest from the root on such a circle would have to choose both incident circle edges, since both would go to nearer vertices. Thus they form a spanning tree $T$.

Every compact subset of this locally finite unit-edge graph lies in a finite subgraph. The open unit balls centred at vertices cover it, a finite subcollection suffices, and each such ball meets only finitely many edges. Given such a finite subgraph and any finitely many specified edges, add the finitely many tree paths from their vertices to the root. The result is a finite connected graph $Y$ with spanning tree $T\cap Y$.

A loop has compact image, so it lies in such a $Y$ and is a word in the loops $\gamma_e$ for edges outside $T$. These loops therefore generate the whole group. If a nonempty reduced word in them contracted in $X_H$, put the compact image of its contraction, together with its edges and root paths, in another such $Y$. The finite graph theorem says those generators are free there, a contradiction. Hence the loops outside $T$ are a free basis.
\end{proof}

\begin{beauty}
Nielsen's 1921 proof used cancellation of words for finitely generated subgroups; Schreier's 1927 proof handled all subgroups and gave the index formula. Coverings explain why a free basis exists: a subgroup becomes a graph, and the edges outside a tree supply the basis. In $F_2$, a subgroup of finite index $k$ has $k+1$ free generators. A subgroup can therefore require more free generators than the group containing it. Infinite index, however, need not mean infinite rank: $\langle a\rangle\le\langle a,b\rangle$ has rank one and infinite index, since its cosets $\langle a\rangle b^j$ are all distinct.
\end{beauty}

\begin{remark}
The theorem also holds for free groups on arbitrary sets. There is a small extra point when vertices have infinitely many incident edges: a compact set can meet infinitely many short edge segments. The finite-subgraph compactness argument above cannot be used unchanged.

Here is a replacement. Subdivide every unit edge into three equal pieces, and cover the graph by open stars of radius $1/5$ about all subdivided vertices. Each star is a contractible tree. Two intersecting stars meet in an interval along their common subdivided edge; there are no triple intersections. A loop can be subdivided so that each piece lies in one star. At a join, move its value to a point on the common edge of the adjacent stars, and join these new values inside each star by its unique finite edge path. The contractions of the stars give a homotopy to a loop using finitely many edges.

The same construction works for a homotopy square after a fine grid subdivision. Assign a star to every rectangle. Move each grid vertex within the intersection of its incident stars, which is a star or an interval, and join the new vertices by finite paths in the intersections appropriate to the edges. Each cell boundary now lies in its assigned star and can be filled by contraction to that star's centre. There are finitely many cells and each contraction uses finitely many edges. If the boundary map already uses finitely many edges, keep it fixed; the same fillings, with its finite boundary pieces, still use finitely many edges. Consequently a finite-edge loop that contracts in the whole graph contracts in a finite subgraph.

Zorn's lemma (Volume~7, Theorem~1.19) supplies a spanning tree: maximize an acyclic edge set containing all vertices; if its forest had two components, a connecting path would allow an extra edge. Every finite collection of edges and its tree paths to a root then sits in a finite connected subgraph whose intersection with the tree spans it. The preceding loop and homotopy reductions let the finite graph theorem prove freeness on all edges outside that tree. Apply this to arbitrary-rank coset graphs to obtain the full Nielsen--Schreier theorem.

Coverings of surfaces give subgroup information too, even though a surface presentation may have relators. As a preview of Chapter~6's Euler-characteristic calculation, a connected finite $k$-sheeted cover of a closed orientable $\Sigma_g$ is a closed orientable $\Sigma_h$ with $2-2h=k(2-2g)$. This finite-cover formula is not being used in the graph proof.
\end{remark}

\section{What we proved, and what comes next}

Going once around a circle leaves a lifted path one unit higher on the line. That integer survives every deformation and adds when loops are joined. On a graph, a spanning tree plays the role of a route back to the starting point; each remaining edge contributes one independent loop. These descriptions give $\Z$ for a circle and free groups for graphs, together with explicit bases for their covering subgroups.

A disc attachment imposes a relation. This turns the bouquet's free group into the surface presentation, and the antipodal sphere covering gives $\pi_1(\R P^n)=\Z/2$ for $n\ge2$. The dimensions matter: $\R P^1$ is a circle, while $\R P^0$ is a point.

Loops still leave some spaces indistinguishable. Both $S^2$ and $S^3$ are simply connected. Chapter~6 builds homology groups to retain cycles of higher dimensions, using abelian groups and boundary maps that can be computed by linear algebra. Those groups will also distinguish Euclidean spaces of different dimensions.

\section{Exercises}

\begin{exercise}[$\star$]\label{ex:fta}
\emph{A polynomial must vanish.} Suppose the monic polynomial $P(z)=z^n+a_{n-1}z^{n-1}+\cdots+a_0$, with $n\ge1$, has no root. Define the loop at $1$
\[
\gamma_r(t)=\frac{P(re^{2\pi it})/P(r)}{|P(re^{2\pi it})/P(r)|},\qquad r\ge0.
\]
(a) Show that its class is independent of $r$. (b) At $r=0$ it is constant. For large $r$, compare $P$ with $z^n$ through a homotopy that keeps the normalized loop based at $1$, and obtain a contradiction.
\end{exercise}

\begin{exercise}[$\star\star$]\label{ex:borsukulam2}
\emph{Antipodal coincidences.} (a) Prove that no map $g:S^2\to S^1$ satisfies $g(-x)=-g(x)$. On the equator, show that such a map would have odd winding number. (b) Deduce that every map $f:S^2\to\R^2$ has $f(x)=f(-x)$ somewhere. Explain the temperature-and-pressure interpretation, assuming continuous measurements on a spherical Earth. (c) Show that $S^2$ cannot be homeomorphic to a subset of $\R^2$.
\end{exercise}

\begin{exercise}[$\star$]\label{ex:figureeight}
\emph{The three two-sheeted covers.} (a) Classify the index-two subgroups of $F_2=\langle a,b\rangle$ as kernels of nonzero maps to $\Z/2$ and identify the covers in Figure~\ref{fig:coverings}. (b) Find a free basis for each subgroup, verify its rank, and decide whether the cover is normal. Justify that the bases in the caption are free, not merely generating sets.
\end{exercise}

\begin{exercise}[$\star\star$]\label{ex:commutator}
\emph{The integer-grid covering.} Let $K=\ker(F_2\to\Z^2)$, where $a,b$ map to the two standard basis vectors. (a) Identify its coset graph with the integer grid, with horizontal $a$-edges and vertical $b$-edges. (b) Exhibit a free basis and prove that its rank is infinite. (c) Show that $K$ is the subgroup generated by all commutators $[g,h]$ in $F_2$.
\end{exercise}

\begin{exercise}[$\star\star$]\label{ex:picture}
\emph{A picture on removable nails.} Let $p\ne q$ be two nails, and model the string as a based loop in $X=\R^2\setminus\{p,q\}$. In this idealized model the picture falls exactly when its loop contracts. (a) Construct a deformation retraction onto a figure eight, with generators $a,b$ encircling the two nails. (b) Removing $q$ erases $b$ and sends $a$ to the generator of $\Z$. Show that $[a,b]$ is nontrivial with both nails but trivial after either is removed. (c) Find a nontrivial loop around three distinct nails that becomes trivial after removal of any one.
\end{exercise}

\begin{exercise}[$\star$]\label{ex:rp3}
\emph{Projective loops and lifted maps.} (a) Prove that every based loop in $\R P^3$ is either nullhomotopic or homotopic to the projection of a half great circle in $S^3$. Show that the second class is nontrivial but its square is trivial. Chapter~7 will identify this loop with one full axial turn in $SO(3)$. (b) Prove that every map from $S^2$ to $S^1$, or from $S^2$ to $T^2$, is homotopic to a constant.
\end{exercise}

\begin{exercise}[$\star\star$]\label{ex:v9-logarithms-roots}
\emph{When can an argument be chosen?} Let $Y$ be path-connected and locally path-connected, and $f:Y\to\C^*$. (a) Prove that a continuous $L:Y\to\C$ with $e^L=f$ exists exactly when every loop $f\gamma$ has winding number zero. (b) For $n\ge1$, prove that a continuous $g$ with $g^n=f$ exists exactly when all these winding numbers are divisible by $n$. If one root exists, show there are exactly $n$. (c) Apply both tests to $f(z)=z^m$ on $1<|z|<2$, for any integer $m$.
\end{exercise}

\begin{exercise}[$\star\star$]\label{ex:v9-three-sheet-cover}
\emph{A cover with no nonidentity deck map.} Use Figure~\ref{fig:v9-three-sheet-cover}, based at vertex $1$. (a) Find the subgroup $H\le F_2$ and a free basis using the tree of $a$-edges $1\to2\to3$. (b) Prove that the deck group is trivial. (c) Check directly that $ab\in H$ but $a^{-1}(ab)a\notin H$, and relate this to normality and the index-three rank formula.
\end{exercise}

\begin{exercise}[$\star\star$]\label{ex:v9-torus-cover-lattice}
\emph{Four sheets can have different symmetries.} Let $m,n>0$ be integers, $r\in\Z$, and $H=\langle(m,0),(r,n)\rangle\le\Z^2$. (a) Construct the covering $\R^2/H\to\R^2/\Z^2$ and show that its total space is a torus and it has $mn$ sheets. (b) Prove that its deck group is
\[
\Z/d_1\times\Z/d_2,\qquad d_1=\gcd(m,r,n),\quad d_2=mn/d_1.
\]
Supply the integer row-and-column reduction needed for this claim. (c) Compare $(m,r,n)=(2,0,2)$ and $(2,1,2)$.
\end{exercise}

\begin{exercise}[$\star\star\star$]\label{ex:v9-shrinking-circles}
\emph{Circles too small to escape a neighbourhood.} Let $C_j$ have centre $(1/j,0)$ and radius $1/j$, and let $E=\bigcup_{j\ge1}C_j\subset\R^2$. (a) Show that $E$ is compact, path-connected and locally path-connected. (b) Prove that the map equal to the identity on $C_j$ and equal to $0$ on all other circles is a continuous retraction $E\to C_j$. (c) Show that $E$ is not semilocally simply connected and has no simply connected covering. Explain why assigning length one to every edge in a graph avoids this obstruction.
\end{exercise}

\section{Hints}

\begin{hint}{ex:fta}
Vary $r$ for the first homotopy. For a large fixed $r$, use $P_s(z)=z^n+s(P(z)-z^n)$, divide $P_s(re^{2\pi it})$ by $P_s(r)$, and normalize its modulus.
\end{hint}
\begin{hint}{ex:borsukulam2}
Lift $h(t)=g(\cos\pi t,\sin\pi t,0)$ for $0\le t\le2$. The difference between the lifted arguments at $t+1$ and $t$ lies in $1/2+\Z$ and is constant. For (b), normalize $f(x)-f(-x)$ if it never vanishes.
\end{hint}
\begin{hint}{ex:figureeight}
An index-two subgroup is normal. In each coset graph choose one edge joining the vertices as a tree. For the third caption basis, give an invertible replacement of one tree-basis element by its product with another.
\end{hint}
\begin{hint}{ex:commutator}
Use the horizontal axis and every vertical line as a spanning tree. The edges outside it are the horizontal edges at heights $j\ne0$. Conjugating a commutator gives another commutator.
\end{hint}
\begin{hint}{ex:picture}
After an affine change put the nails at $(-1,0),(1,0)$. Use two touching diamond boundaries $|x\pm1|+|y|=1$: first push their punctured interiors outwards, then push the exterior vertically onto the diamonds and the horizontal tails. For three nails try $[[a,b],c]$.
\end{hint}
\begin{hint}{ex:rp3}
A lift of a projective loop ends at its starting point or its antipode. For (b), lift the map to $\R$ or $\R^2$ and contract there.
\end{hint}
\begin{hint}{ex:v9-logarithms-roots}
Use the coverings $\exp:\C\to\C^*$ and $z\mapsto z^n$ of $\C^*$. The first has trivial total-space fundamental group; the second multiplies winding by $n$. Two roots have a ratio in the finite set of $n$th roots of unity.
\end{hint}
\begin{hint}{ex:v9-three-sheet-cover}
Tree paths from $1$ are $1,a,a^2$. Read the four edges outside the tree. A deck map must commute with both the cyclic permutation and the transposition of the vertices.
\end{hint}
\begin{hint}{ex:v9-torus-cover-lattice}
Reduce the second coordinate modulo $n$, then the first modulo $m$. For the deck group, use integer row and column operations on $\bigl(\begin{smallmatrix}m&r\\0&n\end{smallmatrix}\bigr)$. Euclidean division reduces the smallest nonzero pivot until it divides every entry.
\end{hint}
\begin{hint}{ex:v9-shrinking-circles}
All circles with sufficiently large index lie in any ball about $0$. Away from $0$, only finitely many circles need be considered. Retract onto one small circle to prove that its loop cannot contract in $E$.
\end{hint}

\section{Solutions}

\begin{solution}{ex:fta}
\idea A root-free polynomial would let a loop of winding number $n$ shrink through the images of smaller circles. Normalize each stage at its starting value so that the basepoint stays fixed.
\writeup (a) The assumption makes $P(r)$ and every numerator nonzero. The formula for $\gamma_r$ depends continuously on $(t,r)$ and equals $1$ at $t=0,1$. Varying $r$ on any closed interval therefore gives a based homotopy. At $r=0$ the ratio is identically $1$.

(b) Choose $r>1+\sum_{j<n}|a_j|$. On $|z|=r$,
\[
|P(z)-z^n|\le r^{n-1}\sum_{j<n}|a_j|<r^n.
\]
Every $P_s(z)=z^n+s(P(z)-z^n)$, $0\le s\le1$, is consequently nonzero on that circle, including at $z=r$. Set
\[
K(t,s)=\frac{P_s(re^{2\pi it})/P_s(r)}{|P_s(re^{2\pi it})/P_s(r)|}.
\]
This is a homotopy of loops based at $1$, joining $\omega_n$ to $\gamma_r$. Their classes are equal, but (a) makes the latter trivial, whereas $[\omega_n]=n\ne0$. The root-free assumption is impossible.
\end{solution}

\begin{solution}{ex:borsukulam2}
\idea Oddness forces an odd number of turns around the equator. A map on the whole sphere would have to make that same loop contract.
\writeup (a) Put $h(t)=g(\cos\pi t,\sin\pi t,0)$ on $[0,2]$, and lift it as $h(t)=e^{2\pi i\alpha(t)}$, choosing any initial argument. For $0\le t\le1$, oddness gives
$\alpha(t+1)-\alpha(t)\in1/2+\Z$. This continuous function into a discrete set is constant, say $q+1/2$. Thus $\alpha(2)-\alpha(0)=2q+1$. On the other hand the equator contracts in $S^2$ by Corollary~\ref{cor:spheres}, so its image has winding number zero. This is a contradiction.

(b) If there were no antipodal coincidence, then
\[
g(x)=\frac{f(x)-f(-x)}{|f(x)-f(-x)|}
\]
would be a continuous odd map to $S^1$, forbidden by (a). Taking the two coordinates to be continuous temperature and pressure gives the stated idealized interpretation. (c) Every map into $\R^2$ identifies an antipodal pair, so none can be injective. In particular there is no homeomorphism onto a subset.
\end{solution}

\begin{solution}{ex:figureeight}
\idea Each generator either fixes or exchanges the two vertices. A connected two-sheeted cover must have at least one exchange.
\writeup (a) An index-two subgroup is normal: its left and right cosets are the subgroup and its complement. Its quotient is $\Z/2$. The possible quotient maps have $(a,b)$ images $(0,1),(1,0),(1,1)$. Conversely every such nonzero map is onto and has an index-two kernel. Distinct maps have distinct kernels, since $\Z/2$ has no nonidentity automorphism. The generator with image $1$ exchanges the sheets, identifying the three drawings.

(b) For $(0,1)$ choose the $b$-edge from the base vertex as the tree. The off-tree edges give $a,bab^{-1},b^2$, a free basis by Theorem~\ref{thm:wedge}. For $(1,0)$ exchange $a,b$, giving $b,aba^{-1},a^2$. For $(1,1)$ choose an $a$-edge as the tree. Its off-tree basis is $a^2,ab,ba^{-1}$. Replace the third element by its product with the second, $b^2$. This is an invertible free-basis change: if the new elements are $A,B,C$, the old third is $CB^{-1}$, while the other two are unchanged. The two substitutions extend to inverse homomorphisms by the free universal property. Hence the caption's set $a^2,ab,b^2$ is also a free basis. Each rank is $3=1+2(2-1)$, and every cover is normal, with the sheet exchange as its nonidentity deck map.
\end{solution}

\begin{solution}{ex:commutator}
\idea A vertex records the two exponent sums of a word. Forgetting everything else turns the covering into the integer grid.
\writeup (a) The map $F_2\to\Z^2$ is onto, so $Kg\mapsto\pi(g)$ identifies cosets with grid vertices. Right multiplication by $a$ or $b$ moves one unit horizontally or vertically.

(b) Let $T$ consist of the horizontal axis and all vertical edges. It is connected and has no circle, since different vertical lines connect only through the axis. The off-tree edges are $(i,j)\to(i+1,j)$ with $i\in\Z$ and $j\in\Z\setminus\{0\}$. Their root-path loops read
\[
a^ib^ja\,b^{-j}a^{-i-1}.
\]
The infinite-graph proof of Nielsen--Schreier makes these a free basis. There are countably infinitely many, proving infinite rank. Infinite index by itself would not suffice for this conclusion.

(c) Every commutator maps to zero in the abelian group $\Z^2$, so the subgroup $C$ they generate lies in $K$. It is normal because $u[g,h]u^{-1}=[ugu^{-1},uhu^{-1}]$. In $F_2/C$ all elements commute, so every class is $a^mb^nC$. The map to $\Z^2$ shows that if the original word is in $K$, then $m=n=0$. Hence $K\subset C$ as well.
\end{solution}

\begin{solution}{ex:picture}
\idea Two independently removable nails call for a word that uses both letters but cancels when either letter is erased. A commutator does exactly that.
\writeup (a) An invertible affine change places $p,q$ at $(-1,0),(1,0)$. Use the two open diamonds $|x+1|+|y|<1$ and $|x-1|+|y|<1$. First push each punctured diamond interior radially onto its boundary. In coordinates $v$ about its centre, with $r=|v_1|+|v_2|\in(0,1)$, the homotopy is $v\mapsto((1-s)+s/r)v$; outside the interiors it is the identity. The formulas agree on their boundaries, including their shared point $0$.

After this step the remaining points satisfy $|y|\ge h(x)$, where
\[
h(x)=\max\{0,1-|x+1|,1-|x-1|\}.
\]
Move such a point vertically by $y\mapsto(1-s)y+s\operatorname{sgn}(y)h(x)$. For $y=0$ the target is zero; there $h(x)=0$, which also ensures continuity as the sign changes. The deformation stays outside the diamond interiors. Its image consists of their two boundaries and the horizontal tails outside $[-2,2]$. Contract each tail towards its endpoint, fixing the diamonds. These successive homotopies fix the two touching diamond boundaries throughout. Their union is a figure eight, so this is a deformation retraction. The two oriented diamond loops freely generate its group.

(b) The word $aba^{-1}b^{-1}$ is reduced and nonempty, hence nontrivial. Erasing $b$ leaves $aa^{-1}$; erasing $a$ leaves $bb^{-1}$. Both contract.

(c) The plane with three punctures has group $F_3$: after rotating so the punctures have different horizontal coordinates, split it into overlapping open half-planes separating one puncture from the other two. The overlap is a puncture-free strip, hence contractible; each half-plane is homeomorphic to a plane with its respective punctures. Van Kampen and (a) give $\Z\ast F_2$. In that group use $[[a,b],c]$. Written out it has ten letters and no adjacent cancellation, so it is nontrivial. Erasing $c$ gives $[a,b][a,b]^{-1}$. Erasing $a$ or $b$ makes the inner commutator trivial and leaves $cc^{-1}$. Each becomes trivial. This is a topological string model; friction and other physical constraints are outside its assumptions. The cited picture-hanging paper develops further constructions.
\end{solution}

\begin{solution}{ex:rp3}
\idea The two possible lifted endpoints give the two loop classes. For maps from a simply connected source, a lift to Euclidean space lets straight segments provide the contraction.
\writeup (a) By Corollary~\ref{cor:spheres}, $S^3\to\R P^3$ is universal, with deck group $\{\mathrm{id},-\mathrm{id}\}$. Thus its base fundamental group is $\Z/2$. A loop is trivial precisely when its lift from $x$ ends at $x$ rather than $-x$. The projection of a half great circle ends upstairs at $-x$, so represents the nontrivial class. Its square represents zero in $\Z/2$. Every nontrivial loop has that same class.

(b) A map $f:S^2\to S^1$ lifts through $\R\to S^1$ because $\pi_1(S^2)$ is trivial. Choose any point over $f(x_0)$ as the initial lift value. Contract the resulting $\tilde f$ by $(1-s)\tilde f$ in $\R$, then project. This joins $f$ to a constant map. The identical argument in $\R^2$ applies to $S^2\to T^2$.
\end{solution}

\begin{solution}{ex:v9-logarithms-roots}
\idea A logarithm or root is a lift through an exponential or power covering. The obstruction is exactly the image of loops in the covering's subgroup.
\writeup (a) On a sufficiently small disc in $\C^*$ choose a continuous argument. Its logarithm branches, differing by $2\pi i k$, show that $\exp:\C\to\C^*$ is a covering. Its total space is convex, so its subgroup is zero. Radial retraction identifies $\pi_1(\C^*)$ with $\Z$ by winding. Choose any logarithm of $f(y_0)$. The lifting criterion gives a lift $L$ with $e^L=f$ exactly when all loop windings vanish.

(b) The same local argument branches show that $p_n(z)=z^n$ covers $\C^*$ with $n$ sheets. A loop of winding $k$ upstairs projects to winding $nk$, so its subgroup is $n\Z$. The criterion gives the divisibility test. If $g$ is one root, all $\zeta g$ with $\zeta^n=1$ are roots. Conversely another root $h$ has continuous ratio $h/g$ in that finite set. Connectedness of $Y$ makes this ratio constant, giving exactly $n$ roots, including one for $n=1$.

(c) The annulus retracts radially onto a circle of radius $3/2$. Its positive generator maps under $z^m$ to winding $m$, also for negative $m$ and for $m=0$. A logarithm therefore exists exactly when $m=0$, and an $n$th root exactly when $n\mid m$. In the latter case the roots are $\zeta z^{m/n}$, $\zeta^n=1$.
\end{solution}

\begin{solution}{ex:v9-three-sheet-cover}
\idea Three sheets need not have three symmetries. A symmetry must preserve the effect of every label, not just permute the vertices arbitrarily.
\writeup (a) Let $H$ consist of the words whose labelled walks start and end at $1$. This is the covering subgroup and has index three. The tree paths to vertices $1,2,3$ read $1,a,a^2$. The off-tree $a$-edge $3\to1$ gives $a^3$. The three $b$-edges $1\to2$, $2\to1$, $3\to3$ give, respectively,
\[
ba^{-1},\qquad ab,\qquad a^2ba^{-2}.
\]
Together with $a^3$ these are a free basis by the graph theorem. The rank is $6-3+1=4$, agreeing with $1+3(2-1)$.

(b) A deck map's permutation of the three vertices commutes with the $a$-cycle. Such a permutation is determined by the image of $1$ and must be one of its three powers. The two nonidentity powers do not commute with the $b$-transposition: for example they carry its fixed vertex $3$ elsewhere. Thus only the identity permutation is possible. A deck map fixing a vertex is the identity on the whole connected cover, so the deck group is trivial.

(c) Reading $ab$ goes $1\to2\to1$, so $ab\in H$. But $a^{-1}(ab)a=ba$, which goes $1\to2\to3$, so is not in $H$. Therefore $H$ is not normal and the cover is not normal. The deck theorem further gives $N(H)=H$, consistent with the trivial deck group.
\end{solution}

\begin{solution}{ex:v9-torus-cover-lattice}
\idea The columns of an integer matrix describe which translations have already been identified upstairs. Its determinant counts sheets; its integer reduction determines the remaining translation symmetries.
\writeup (a) Let $A=\bigl(\begin{smallmatrix}m&r\\0&n\end{smallmatrix}\bigr)$. The invertible linear map $A:\R^2\to\R^2$ sends $\Z^2$ onto $H$, inducing a homeomorphism $T^2\cong\R^2/H$. Inclusion $H\subset\Z^2$ gives the stated projection. Small discs of radius less than $1/2$ in $\R^2$ have disjoint nontrivial integer translates; their images give local sheets for the projection. Its induced subgroup is $H$, since a loop closes upstairs precisely when a lifted path in $\R^2$ ends at a vector of $H$.

Every integer vector has a unique representative modulo $H$ with second coordinate in $\{0,\ldots,n-1\}$ and first coordinate in $\{0,\ldots,m-1\}$: subtract a multiple of $(r,n)$ first, then of $(m,0)$. Uniqueness follows by comparing the second, then first coordinates. There are $mn$ cosets, hence $mn$ sheets. As $\Z^2$ is abelian, the covering is normal with deck group $\Z^2/H$.

(b) Integer row operations change the ambient basis of $\Z^2$, and integer column operations change the generators of $H$; swaps, sign changes and adding integer multiples are invertible. They therefore preserve the quotient up to isomorphism. For completeness, reduce a nonsingular $2\times2$ integer matrix as follows. Move a nonzero entry to the upper left and make it positive, calling it $d$. Euclidean row and column division either reduces $d$ to a strictly smaller positive remainder or clears its column and row, leaving $\operatorname{diag}(d,c)$. If $d\nmid c$, add the second row to the first; the first row now contains $d,c$, and Euclidean column division produces a positive remainder smaller than $d$. Repeat. The positive pivot cannot decrease indefinitely, so eventually we obtain $\operatorname{diag}(d,c)$ with $d\mid c$; change the final sign if needed.

The greatest common divisor of all entries is unchanged by every operation. In this final matrix it is $d$, so $d=d_1=\gcd(m,r,n)$. The absolute determinant is also unchanged, so $dc=mn$ and $c=d_2$. The quotient of this diagonal subgroup is $\Z/d_1\times\Z/d_2$, as claimed.

(c) For $(2,0,2)$ the group is $\Z/2\times\Z/2$. For $(2,1,2)$ it is $\Z/4$. Both are four-sheeted torus covers, but their deck groups are not isomorphic, so the coverings cannot be isomorphic.
\end{solution}

\begin{solution}{ex:v9-shrinking-circles}
\idea A small circle is still a whole circle. Retraction onto it detects a loop that persists inside every neighbourhood of the common point.
\writeup (a) The set is bounded. A convergent sequence with points on circles of unbounded index has a subsequence tending to $0$, because $C_j$ lies in the ball of radius $2/j$. A sequence with bounded indices has a subsequence in one of finitely many closed circles. This proves that every limit point lies in $E$, so $E$ is closed and compact. Each point joins $0$ along its circle, proving path-connectedness.

Away from $0$, all sufficiently small circles are absent from a sufficiently small neighbourhood. The remaining finitely many circles meet one another only at $0$, so a smaller neighbourhood of the given point is an arc. At $0$, the sets $E\cap B_\varepsilon(0)$ form a path-connected neighbourhood basis: each intersection with a circle is either that whole circle or the arc through $0$ cut out by the ball. Thus $E$ is locally path-connected.

(b) Let $r_j$ be the stated map. At $0$, $|r_j(x)|\le|x|$, giving continuity. At a nonzero point of $C_j$, choose a small neighbourhood meeting none of the other circles, as in (a); there $r_j$ is the identity. At a point outside the closed set $C_j$, a neighbourhood misses $C_j$, and $r_j$ is constant. Hence it is a continuous retraction.

(c) Every neighbourhood of $0$ contains $C_j$ for all sufficiently large $j$. A loop going once around such a circle cannot contract in $E$, because composing a contraction with $r_j$ would contract its generator in $C_j$, whose group is $\Z$. The semilocal condition therefore fails.

Suppose nevertheless that $p:\tilde E\to E$ were a covering with simply connected total space. Choose an evenly covered neighbourhood $U$ of $0$ and a sheet containing a point over $0$. A sufficiently small $C_j$ lies in $U$, so its generating loop lifts by the inverse branch to a closed loop in that sheet. It contracts in $\tilde E$; projecting would contract it in $E$, a contradiction. In a unit-edge graph, by contrast, a ball of radius less than $1/2$ at a vertex is a star with no whole circular edge inside it, and contracts along its arms. Shrinking edge lengths destroy that uniform local description.
\end{solution}

\section*{Sources and further reading}
\addcontentsline{toc}{section}{Sources and further reading}

H.~Poincar\'e, ``Analysis situs'', \emph{J.\ \'Ec.\ Polytech.}\ (2) 1 (1895), 1--121, and ``Cinqui\`eme compl\'ement \`a l'Analysis situs'', \emph{Rend.\ Circ.\ Mat.\ Palermo} 18 (1904), 45--110; both are
translated in H.~Poincar\'e, \emph{Papers on Topology: Analysis Situs and Its Five Supplements}, translated by J.~Stillwell (AMS and LMS, 2010), whose introduction tells the story of the corrections. On the
solution: G.~Perelman, arXiv preprints math/0211159, math/0303109 and math/0307245 (2002--2003); J.~Morgan and G.~Tian, \emph{Ricci Flow and the Poincar\'e Conjecture} (AMS and Clay Mathematics
Institute, 2007); the Clay Mathematics Institute's announcement of 18~March~2010; J.~Milnor, ``Towards the Poincar\'e conjecture and the classification of 3-manifolds'', \emph{Notices AMS} 50 (2003), 1226--1233.

H.~Seifert, ``Kon\-struk\-tion drei\-di\-men\-sio\-na\-ler ge\-schlos\-se\-ner R\"aume'', \emph{Ber.\ S\"achs.\ Akad.\ Wiss.} 83 (1931), 26--66; E.~R.~van Kampen, ``On the connection between the fundamental groups of
some related spaces'', \emph{Amer.\ J.\ Math.} 55 (1933), 261--267. J.~Nielsen, ``Om regning med ikke-kommutative faktorer og dens anvendelse i gruppeteorien'', \emph{Mat.\ Tidsskr.\ B} (1921), 77--94;
O.~Schreier, ``Die Untergruppen der freien Gruppen'', \emph{Abh.\ Math.\ Sem.\ Hamburg} 5 (1927), 161--183. E.~D.~Demaine, M.~L.~Demaine, Y.~N.~Minsky, J.~S.~B.~Mitchell, R.~L.~Rivest and M.~P\v atra\c scu,
``Picture-hanging puzzles'', \emph{Theory Comput.\ Syst.} 54 (2014), 531--550.

For a second route through the lifting and gluing arguments, read A.~Hatcher, \emph{Algebraic Topology} (Cambridge University Press, 2002), Sections~1.1--1.3 and~1.A. Hatcher also treats van Kampen for more than two open sets. W.~Massey, \emph{Algebraic Topology: An Introduction} (Harcourt, Brace and World, 1967), develops the geometric viewpoint through many examples. The permutation proof of Proposition~\ref{prop:freeproduct} comes from van der Waerden, ``Free products of groups'', \emph{Amer.\ J.\ Math.} 70 (1948), 527--528. To see how finite graphs become tools for studying subgroups, continue with J.~Stallings, ``Topology of finite graphs'', \emph{Invent.\ Math.} 71 (1983), 551--565.
